Bayesian optimization (BO) is a class of methods for global optimization of expensive black-box functions. 
In BO, a probabilistic surrogate model is used to approximate the objective and future evaluations are selected via an acquisition function. 
BO has been successfully applied to applications such as robotic gait control~\citep{calandra2016bayesian}, sensor set selection~\citep{garnett2010bayesian}, and neural network hyperparameter tuning~\citep{snoek2012practical,turner2021bayesian}.
BO is favored for these tasks because of its sample-efficient nature.
Achieving this sample-efficiency requires BO to balance exploration and exploitation. 
However, standard acquisition functions such as expected improvement (EI) are often too greedy in practice. 
As a result, they perform poorly on multimodal problems~\citep{hernandez2014predictive} and have provably sub-optimal performance in certain settings, e.g., bandit problems~\citep{srinivas2009gaussian}. 
A key research goal in BO is developing less greedy acquisition functions~\citep{shahriari2016botutorial}. 
Examples include predictive entropy search (PES)~\citep{hernandez2014predictive} or knowledge gradient (KG)~\citep{frazier2008knowledge}. 
\citet{lam2016bayesian} frame the exploration-exploitation trade-off as a balance between immediate and future rewards in a continuous state and action space Markov decision process (MDP). 
In this framework, \textit{non-myopic} acquisition functions are optimal MDP policies, and promise better performance by considering the impact of future evaluations up to a given BO budget (also referred to as the \textit{horizon}). 

While BO budgets are typically given in iterations, this implicitly measures convergence in terms of iteration count and assumes uniform evaluation cost. 
For many practical BO applications, evaluation costs may vary in different regions of the search space. 
For example, the time spent on neural network training increases with layer size; the cost of different drug compounds vary; and control actions in optimal control have differing complexities. 
In all these cases, standard BO is often unable to achieve fast convergence in terms of unit cost. 
Motivated by these examples, we develop methods that improve convergence when measured by cost. 
This cost may be time, energy, or money, and the goal is to minimize the objective given a cost budget. 

\begin{figure*}[!ht] 
    \centering
    \includegraphics[width=\textwidth]{histogram_variations.pdf}
    \caption{Runtime distribution, log-scaled, of $5000$ randomly selected points for the $k$-nearest-neighbors (KNN), Multi-layer Perceptron (MLP), Support Vector Machine (SVM), Decision Tree (DT), and Random Forest (RF) hyperparameter optimization problems, each trained on the OpenML w2a dataset ~\citep{OpenML2013}. The runtimes vary, often by an order of magnitude or more.}
    \label{fig:variations}
\end{figure*}

\textit{Cost-constrained BO} measures convergence with these alternative cost metrics for which standard BO methods are unsuited. 
We extend non-myopic BO to handle the cost-constrained setting. Our contributions follow:
\begin{itemize}
    \item We analyze failure modes of common approaches to cost-constrained BO, in which greedy behavior results in poor per-cost performance.  
    \item To avoid overly greedy behavior, we formulate cost-constrained BO as an instance of a constrained Markov decision process (CMDP). This formulation is a novel extension of recent research on non-myopic BO which uses a simpler Markov decision process (MDP) formulation. 
    \item We introduce an approximation to the optimal constrained MDP policy based on rollout of feasible trajectories. Rollout is a popular class of approximate MDP solutions in which future BO realizations and their corresponding values are simulated using the surrogate and then averaged.
    \item We validate the performance of our methods on a set of practical hyperparameter optimization problems and a sensor set selection problem.
\end{itemize}
